\SetPageTheme{story}
\PageHeader
  {Structura pentru}
  {Cand intervalul este clar de la inceput, putem strange numararea intr-o forma mai scurta si mai ordonata.}
  {Pentru 1 din 6}

\begin{missionbox}[De ce mai avem nevoie de ea]
Structura \texttt{cat timp} ne invata foarte bine ideea de repetitie. Dar multe probleme sunt deja planificate:
\begin{itemize}
\item stii de la ce valoare pornesti;
\item stii pana unde mergi;
\item stii cu ce pas avansezi.
\end{itemize}

In aceste situatii, structura \texttt{pentru} este mai lizibila si mai usor de verificat.
\end{missionbox}

\BlocklyExperiment{cq-i1-for-interval}{Parcurgere pe interval}{Modifica startul, limita si pasul, apoi compara secventele obtinute si acoperirea intervalului.}

\clearpage

\SetPageTheme{theory}
\PageHeader
  {Cele trei componente}
  {Startul, limita si pasul trebuie gandite impreuna, nu separat.}
  {Pentru 2 din 6}

\begin{examplebox}[Modelul de baza]
\begin{lstlisting}[style=codequest]
pentru i de la 1 la 5 cu pas 1 executa
    afiseaza i
sfarsit
\end{lstlisting}
\end{examplebox}

\begin{conceptbox}[Ce spune dintr-o privire]
\begin{itemize}
\item \textbf{de la 1} inseamna valoarea de start;
\item \textbf{la 5} inseamna limita;
\item \textbf{cu pas 1} inseamna ritmul de parcurgere.
\end{itemize}
Aceasta claritate este motivul pentru care \texttt{pentru} este atat de folosita in randare si in parcurgerea grilelor.
\end{conceptbox}

\begin{guidebox}[Identificare]
Pentru fiecare secventa, noteaza startul, limita si pasul:
\begin{enumerate}
\item \texttt{pentru i de la 1 la 10 cu pas 1}
\item \texttt{pentru i de la 2 la 20 cu pas 2}
\item \texttt{pentru i de la 10 la 0 cu pas -1}
\end{enumerate}
\AnswerSpace{7}
\end{guidebox}

\clearpage

\SetPageTheme{guided}
\PageHeader
  {Filtram in interiorul intervalului}
  {Parcurgerea si selectia merg foarte des impreuna.}
  {Pentru 3 din 6}

\begin{examplebox}[Pare dintr-un interval]
\begin{lstlisting}[style=codequest]
pentru i de la 1 la 12 cu pas 1 executa
    daca (i este par) atunci
        afiseaza i
    sfarsit
sfarsit
\end{lstlisting}
\end{examplebox}

\begin{missionbox}[Legatura cu jocurile]
Exact asa functioneaza multe mini-sarcini din jocuri:
\begin{itemize}
\item verifici toate celulele de pe un rand;
\item verifici toate monedele dintr-o lista;
\item verifici toate pozitiile de pe o linie de harta;
\item marchezi doar elementele care respecta o regula.
\end{itemize}
\end{missionbox}

\begin{semibox}[Aplicam]
Scrie doua variante:
\begin{enumerate}
\item afiseaza multiplii lui 3 pana la 24;
\item afiseaza doar valorile impare dintre 1 si 15.
\end{enumerate}
\AnswerSpace{8}
\end{semibox}

\clearpage

\SetPageTheme{guided}
\PageHeader
  {Podul catre randare}
  {Acelasi tipar de interval devine, mai tarziu, o linie de celule sau o linie de pixeli.}
  {Pentru 4 din 6}

\begin{missionbox}[Cum se leaga de partea vizuala]
Cand randam o linie, programul parcurge pozitii una cate una:
\begin{center}
\textit{de la stanga la dreapta, de la coloana 1 la coloana n}
\end{center}
De aceea structura \texttt{pentru} este atat de potrivita pentru tabla de joc, matrice si desenarea formelor.
\end{missionbox}

\begin{guidebox}[Aplicatie pe o linie de tabla]
Construieste o bucla care parcurge coloanele 1..8 si afiseaza:
\begin{itemize}
\item \texttt{X} pe coloanele pare;
\item \texttt{.} pe coloanele impare.
\end{itemize}

Scrie si o propozitie scurta despre cum ai traduce acelasi model pentru un rand complet dintr-un joc.
\AnswerSpace{10}
\end{guidebox}

\clearpage

\SetPageTheme{simulation}
\PageHeader
  {Simularea intervalului}
  {Aici vezi imediat ce castigi si ce pierzi atunci cand schimbi pasul.}
  {Pentru 5 din 6}

\SimulatorCard{cq-i1-for-interval}{Interval si pas}{Felul in care pasul modifica numarul de iteratii, valorile atinse si densitatea parcurgerii.}{Testeaza acelasi interval cu pasi diferiti, apoi explica ce valori au fost vizitate si ce valori au fost sarite.}

\begin{semibox}[Dupa scanare]
\begin{enumerate}
\item Cand pasul este 1, ce avantaj de claritate ai?
\item Cand pasul este mare, ce risc apare?
\item Ce pas ai alege pentru a desena fiecare celula de pe un rand? Dar pentru a vizita doar anumite coloane?
\end{enumerate}
\AnswerSpace{6}
\end{semibox}

\clearpage

\SetPageTheme{challenge}
\PageHeader
  {Scrii intervale complete}
  {Acum alegi singur forma potrivita si explici logica din spatele ei.}
  {Pentru 6 din 6}

\SyntaxCheck{pentru}

\begin{solobox}[Set complet]
\begin{enumerate}
\item Scrie o bucla pentru toate valorile 1..30 divizibile cu 5.
\item Scrie o bucla descrescatoare 21..0 cu pas 3.
\item Scrie o bucla care afiseaza doar valorile ce trec o conditie aleasa de tine.
\item Explica in 4--5 randuri de ce structura \texttt{pentru} este mai potrivita decat \texttt{cat timp} atunci cand intervalul este deja cunoscut.
\end{enumerate}
\AnswerSpace{12}
\end{solobox}

\begin{checkpointbox}[Ideea care ramane]
\texttt{Pentru} este forma cea mai buna cand numararea este planificata: interval cunoscut, pas clar, rezultat usor de verificat.
\end{checkpointbox}

\clearpage
